\begin{frame}{Conclusiones}
  \begin{columns}[T]
    \begin{column}{.48\textwidth}
      \begin{block}{Sobre el artículo}
        \footnotesize
        \begin{itemize}
          \item La idea es fuerte y simple: una CNN con parches y diseño isotrópico compite con los transformers.
          \item La evidencia es parcial: sin diseño factorial, con recetas distintas entre modelos y una sola corrida
          por configuración.
          \item Sí muestra algo valioso: la receta de entrenamiento pesa tanto como la arquitectura.
        \end{itemize}
      \end{block}
    \end{column}
    \begin{column}{.48\textwidth}
      \begin{block}{Sobre la réplica}
        \footnotesize
        \begin{itemize}
          \item Con la receta oficial en GPU: \accColab\,\% (paper: 95,21\,\%). La implementación es correcta.
          \item Con poco cómputo: ConvMixer $\approx$ ResNet (\accCMnueve{} vs.\ \accResNet\,\%), pero la ResNet es
          \ratioTput{} veces más rápida.
          \item Se reproduce el efecto del parche; el del núcleo grande, no con 10 épocas.
          \item Modelos del curso: lineal $\approx$ 40 $<$ ELM 47 $<$ MLP 57 $\ll$ CNN 83--85\,\%.
        \end{itemize}
      \end{block}
    \end{column}
  \end{columns}
  \clave{Replicar no es copiar un número: es seguir el procedimiento, medir lo mismo y poder explicar cada
  diferencia. Aquí cada brecha tiene una causa concreta: épocas, tamaño, aumentación o la semilla.}
\end{frame}
